%!TEX program = xelatex
\documentclass[12pt, a4paper]{ctexart}
\usepackage[normalem]{ulem}
\usepackage{enumitem}
\usepackage{amsmath}
\usepackage{amssymb}
\usepackage{tcolorbox}
\tcbuselibrary{breakable}
\usepackage{url}
\usepackage{geometry}
\usepackage{booktabs}
\usepackage{longtable}
\usepackage{xcolor}
\usepackage{listings}
\lstset{
    basicstyle=\ttfamily\small,
    breaklines=true,
    frame=single,
    backgroundcolor=\color{gray!10},
    numbers=left,
    numberstyle=\tiny\color{gray},
}
\geometry{left=2.5cm, right=2.5cm, top=2.5cm, bottom=2.5cm}
\setcounter{secnumdepth}{3}

\setlistdepth{5}
\renewlist{itemize}{itemize}{5}
\renewlist{enumerate}{enumerate}{5}
\setlist[itemize,1]{label=\textbullet}
\setlist[itemize,2]{label=--}
\setlist[itemize,3]{label=*}
\setlist[enumerate,1]{label=\arabic*.}
\setlist[enumerate,2]{label=\alph).}
\setlist[enumerate,3]{label=\roman).}

\newenvironment{note}{%
  \begin{tcolorbox}[colback=blue!5, colframe=blue!50!black, boxrule=0.5pt, arc=2pt, breakable, title=\textbf{说明}]
}{%
  \end{tcolorbox}
}

\newenvironment{warning}{%
  \begin{tcolorbox}[colback=red!5, colframe=red!50!black, boxrule=0.5pt, arc=2pt, breakable, title=\textbf{注意}]
}{%
  \end{tcolorbox}
}

\newenvironment{codefile}{%
  \begin{tcolorbox}[colback=gray!5, colframe=gray!50!black, boxrule=0.5pt, arc=0pt, breakable]
}{%
  \end{tcolorbox}
}

\title{\textbf{LDA技能演化分析操作手册} \\ \large 基于招聘数据的动态主题建模}
\author{操作化细则 v3.0}
\date{2026年1月}

\begin{document}

\maketitle
\tableofcontents
\newpage

%==============================================================================
\section{项目概述}
%==============================================================================

\subsection{研究目标}
基于2016--2024年招聘数据（约1300万条），使用LDA主题模型识别职业技能组合，并追踪技能需求的时间演化趋势。

\subsection{数据现状}

\begin{table}[h]
\centering
\caption{数据资源清单}
\begin{tabular}{lll}
\toprule
\textbf{资源} & \textbf{路径} & \textbf{状态} \\
\midrule
清洗后数据 & \texttt{data/processed/cleaned/all\_in\_one1.csv} & 已完成 (13GB) \\
任职要求提取 & \texttt{data/processed/extracted/all\_in\_one2.csv} & 已完成 (5.6GB) \\
去重后数据 & \texttt{data/processed/deduped/all\_in\_one2\_dedup.csv} & 已完成 (4.5GB) \\
ESCO技能表（EN） & \texttt{data/esco/original/ESCO dataset.../skills\_en.csv} & 已完成 \\
ESCO分词词典 & \texttt{data/esco/jieba\_dict/esco\_jieba\_dict.txt} & 已完成（27,773词，中英双语） \\
ESCO技能元数据 & \texttt{data/esco/jieba\_dict/skill\_metadata.csv} & 已完成（含 hard/soft/lang 分类） \\
\midrule
\multicolumn{3}{l}{\textbf{时间窗口数据（Task 1.2 输出）}} \\
Phase I 数据 & \texttt{data/processed/windows/window\_2016\_2018.csv} & 已完成（2,132,073条，23.3\%） \\
Phase II 数据 & \texttt{data/processed/windows/window\_2021\_2022.csv} & 已完成（4,186,033条，45.7\%） \\
Phase III 数据 & \texttt{data/processed/windows/window\_2023\_2024.csv} & 已完成（2,809,953条，30.7\%） \\
\bottomrule
\end{tabular}
\end{table}

\subsection{年份分布与时间窗口设计}

\begin{warning}
数据在2019--2020年存在明显断层（2019年仅18条，2020年仅2054条），无法支持逐年分析。
\end{warning}

\textbf{时间窗口划分方案}：

\begin{table}[h]
\centering
\caption{时间窗口设计（AI技术代际视角）}
\begin{tabular}{llrl}
\toprule
\textbf{窗口} & \textbf{年份} & \textbf{实际记录数} & \textbf{技术背景} \\
\midrule
Phase I & 2016--2018 & 2,132,073 (23.3\%) & AI 1.0: 深度学习渗透期 \\
Phase II & 2021--2022 & 4,186,033 (45.7\%) & 数字化转型深化期 \\
Phase III & 2023--2024 & 2,809,953 (30.7\%) & AI 2.0: 大模型爆发期 \\
\bottomrule
\end{tabular}
\end{table}

\begin{note}
\textbf{窗口划分依据}：基于AI技术代际而非外部事件。2019--2020数据断层恰好对应技术迭代分水岭：Phase I 为AlphaGo后深度学习早期渗透；Phase III 为ChatGPT引发的大模型革命。
\end{note}

%==============================================================================
\section{Phase 1: 词典构建与分词}
%==============================================================================

\subsection{Task 1.1: ESCO词典转jieba格式}

\textbf{目标}：将ESCO技能词表转换为jieba用户词典，防止专业术语被错误切分。

\textbf{输入}：\texttt{data\_cleaning/ESCO dataset - v1.2.1 - classification - en - csv/skills\_en.csv}

\textbf{输出}：
\begin{itemize}
    \item \texttt{Test\_LDA/task1\_datacleaning/esco\_jieba\_dict.txt}
    \item \texttt{data\_cleaning/jieba/skill\_metadata.csv}（含 soft/hard 标签）
    \item \texttt{data\_cleaning/jieba/skills\_ch.csv}
\end{itemize}

\textbf{状态}：已完成（2026年1月28日）

\begin{codefile}
\textbf{文件}: \texttt{data\_cleaning/jieba/generate\_final\_dict.py}
\begin{lstlisting}[language=Python]
#!/usr/bin/env python3
# -*- coding: utf-8 -*-
"""
从清洗后的 ESCO 元数据生成 jieba 词典（所见即所得）
"""
import pandas as pd
import re
from pathlib import Path

# === 配置路径（绝对路径）===
INPUT_META = '/Users/yu/code/code2601/TY/data_cleaning/jieba/skill_metadata.csv'
OUTPUT_DICT = '/Users/yu/code/code2601/TY/Test_LDA/task1_datacleaning/esco_jieba_dict.txt'

def main():
    print(f"Reading source of truth: {INPUT_META}...")
    df = pd.read_csv(INPUT_META)

    if 'term' not in df.columns:
        print("Error: 'term' column missing in metadata!")
        return

    unique_terms = df['term'].dropna().unique()
    print(f"Found {len(unique_terms)} unique terms to strictly enforce.")

    Path(OUTPUT_DICT).parent.mkdir(parents=True, exist_ok=True)
    with open(OUTPUT_DICT, 'w', encoding='utf-8') as f:
        for term in unique_terms:
            term_str = str(term).strip()
            if not term_str:
                continue
            f.write(f"{term_str} 20000 nz\n")

    print(f"Successfully generated {OUTPUT_DICT}.")

if __name__ == '__main__':
    main()
\end{lstlisting}
\end{codefile}

\textbf{执行命令}：
\begin{verbatim}
python /Users/yu/code/code2601/TY/data_cleaning/jieba/generate_final_dict.py
\end{verbatim}

%------------------------------------------------------------------------------
\subsection{Task 1.2: 按时间窗口分割数据}

\textbf{目标}：将去重后的招聘数据按AI技术代际划分为三个时间窗口。

\textbf{输入}：\texttt{data/processed/deduped/all\_in\_one2\_dedup.csv}

\textbf{输出}：\texttt{data/processed/windows/window\_YYYY\_YYYY.csv}

\textbf{状态}：已完成（2026年1月29日）

\textbf{执行结果}：
\begin{itemize}
    \item Phase I (2016-2018): 2,132,073 条 (23.3\%)
    \item Phase II (2021-2022): 4,186,033 条 (45.7\%)
    \item Phase III (2023-2024): 2,809,953 条 (30.7\%)
    \item 未纳入分析: 32,077 条 (2019: 17条, 2020: 1,479条, 2025: 30,581条)
\end{itemize}

\begin{codefile}
\textbf{文件}: \texttt{src/cleaning/split\_by\_window.py}
\begin{lstlisting}[language=Python]
#!/usr/bin/env python3
# -*- coding: utf-8 -*-
"""
按时间窗口分割数据
"""
import pandas as pd
import sys
from pathlib import Path
from tqdm import tqdm

# 路径配置
sys.path.insert(0, str(Path(__file__).parent.parent.parent))
from config.paths import DEDUPED_DATA, PROCESSED_DIR

INPUT_FILE = DEDUPED_DATA  # data/processed/deduped/all_in_one2_dedup.csv
OUTPUT_DIR = PROCESSED_DIR / 'windows'
OUTPUT_DIR.mkdir(parents=True, exist_ok=True)
CHUNKSIZE = 100000

# 时间窗口定义
WINDOWS = {
    'window_2016_2018': [2016, 2017, 2018],
    'window_2021_2022': [2021, 2022],
    'window_2023_2024': [2023, 2024],
}

def main():
    Path(OUTPUT_DIR).mkdir(parents=True, exist_ok=True)

    # 初始化输出文件
    file_handles = {}
    first_chunk = {w: True for w in WINDOWS}
    counts = {w: 0 for w in WINDOWS}

    print("分割数据到各时间窗口...")

    for chunk in tqdm(pd.read_csv(INPUT_FILE, chunksize=CHUNKSIZE)):
        for window_name, years in WINDOWS.items():
            # 筛选对应年份
            mask = chunk['招聘发布年份'].isin(years)
            window_data = chunk[mask]

            if len(window_data) > 0:
                output_file = Path(OUTPUT_DIR) / f"{window_name}.csv"
                window_data.to_csv(
                    output_file,
                    mode='w' if first_chunk[window_name] else 'a',
                    index=False,
                    header=first_chunk[window_name]
                )
                first_chunk[window_name] = False
                counts[window_name] += len(window_data)

    # 输出统计
    print("\n分割完成:")
    for window_name, count in counts.items():
        print(f"  {window_name}: {count:,} 条")

if __name__ == '__main__':
    main()
\end{lstlisting}
\end{codefile}

%------------------------------------------------------------------------------
\subsection{Task 1.3: 分词与短语识别}

\textbf{目标}：使用jieba+ESCO词典分词，并用Phrases模型识别固定搭配。

\textbf{输入}：\texttt{data/processed/windows/window\_*.csv}

\textbf{输出}：\texttt{data/processed/tokenized/}
\begin{itemize}
    \item \texttt{window\_*\_corpus.pkl} -- 分词后语料
    \item \texttt{window\_*\_phraser.pkl} -- Bigram/Trigram短语模型
    \item \texttt{window\_*\_tokenized.csv} -- 带tokenized列的CSV（便于检查）
\end{itemize}

\textbf{状态}：已完成（2026年1月29日）

\textbf{执行结果}：
\begin{itemize}
    \item Phase I (2016-2018): 2,131,984 文档, 词表 313,614词, 平均 28.7词/文档
    \item Phase II (2021-2022): 4,185,948 文档, 词表 379,436词, 平均 28.8词/文档
    \item Phase III (2023-2024): 2,809,898 文档, 词表 313,353词, 平均 31.0词/文档
    \item 总计: 9,127,830 有效文档
\end{itemize}

\begin{codefile}
\textbf{文件}: \texttt{src/cleaning/tokenize\_with\_esco.py}
\begin{lstlisting}[language=Python]
#!/usr/bin/env python3
# -*- coding: utf-8 -*-
"""
使用ESCO词典进行分词+短语识别
"""
import jieba
import pandas as pd
import pickle
from pathlib import Path
from gensim.models import Phrases
from gensim.models.phrases import Phraser
from tqdm import tqdm
import re

# 路径配置
ESCO_DICT = '/Users/yu/code/code2601/TY/Test_LDA/task1_datacleaning/esco_jieba_dict.txt'
WINDOW_DIR = '/Users/yu/code/code2601/TY/Test_LDA/task1_datacleaning/windows'
OUTPUT_DIR = '/Users/yu/code/code2601/TY/Test_LDA/task1_datacleaning/tokenized'

# 加载ESCO词典
print("加载ESCO词典...")
jieba.load_userdict(ESCO_DICT)

# 停用词（招聘文本专用）
STOP_WORDS = set([
    # 通用停用词
    '的', '了', '在', '是', '有', '和', '与', '或', '等', '能', '会', '及',
    '对', '可', '为', '被', '把', '让', '给', '向', '从', '到', '以', '于',
    # 招聘高频无意义词
    '负责', '具有', '具备', '熟悉', '掌握', '了解', '相关', '以上', '以下',
    '优先', '考虑', '经验', '工作', '能力', '要求', '岗位', '公司', '企业',
    '团队', '良好', '较强', '一定', '优秀', '专业', '学历', '年以上', '年',
    '及以上', '左右', '不限', '若干', '可以', '需要', '进行', '开展',
    '完成', '参与', '支持', '提供', '负责', '协助', '配合', '执行',
])

def clean_text(text: str) -> str:
    """清理文本"""
    if pd.isna(text):
        return ""
    text = str(text)
    # 去除特殊字符，保留中英文、数字、常见符号
    text = re.sub(r'[^\u4e00-\u9fa5a-zA-Z0-9\.\+\#\-\_\s]', ' ', text)
    text = re.sub(r'\s+', ' ', text)
    return text.strip()

def tokenize(text: str) -> list:
    """分词并过滤"""
    text = clean_text(text)
    if not text:
        return []

    words = jieba.lcut(text)
    # 过滤：长度>=2，非停用词，非纯数字
    words = [w for w in words
             if len(w) >= 2
             and w not in STOP_WORDS
             and not w.isdigit()
             and not re.match(r'^[\d\.]+$', w)]
    return words

def process_window(window_file: str):
    """处理单个时间窗口"""
    window_name = Path(window_file).stem
    print(f"\n处理 {window_name}...")

    # 读取数据
    df = pd.read_csv(window_file)
    print(f"  记录数: {len(df):,}")

    # Step 1: 基础分词
    print("  Step 1: 分词...")
    corpus = []
    for text in tqdm(df['cleaned_requirements'], desc="分词"):
        tokens = tokenize(text)
        if tokens:  # 只保留非空文档
            corpus.append(tokens)

    print(f"  有效文档: {len(corpus):,}")

    # Step 2: 短语识别（使用部分样本训练）
    print("  Step 2: 训练短语模型...")
    sample_size = min(500000, len(corpus))
    sample = corpus[:sample_size]

    # Bigram
    bigram = Phrases(sample, min_count=100, threshold=15)
    bigram_phraser = Phraser(bigram)

    # Trigram
    trigram = Phrases(bigram_phraser[sample], min_count=50, threshold=15)
    trigram_phraser = Phraser(trigram)

    # Step 3: 应用短语模型
    print("  Step 3: 应用短语模型...")
    corpus = [list(trigram_phraser[bigram_phraser[doc]]) for doc in tqdm(corpus, desc="应用短语")]

    # 保存
    Path(OUTPUT_DIR).mkdir(parents=True, exist_ok=True)

    corpus_file = Path(OUTPUT_DIR) / f"{window_name}_corpus.pkl"
    with open(corpus_file, 'wb') as f:
        pickle.dump(corpus, f)
    print(f"  保存语料: {corpus_file}")

    phraser_file = Path(OUTPUT_DIR) / f"{window_name}_phraser.pkl"
    with open(phraser_file, 'wb') as f:
        pickle.dump((bigram_phraser, trigram_phraser), f)
    print(f"  保存短语模型: {phraser_file}")

    # 保存带tokenized列的CSV（用于后续分析）
    df_valid = df.iloc[:len(corpus)].copy()
    df_valid['tokenized'] = [' '.join(doc) for doc in corpus]
    tokenized_file = Path(OUTPUT_DIR) / f"{window_name}_tokenized.csv"
    df_valid.to_csv(tokenized_file, index=False)

    return len(corpus)

def main():
    total = 0
    for window_file in sorted(Path(WINDOW_DIR).glob('window_*.csv')):
        count = process_window(str(window_file))
        total += count

    print(f"\n全部完成！总文档数: {total:,}")

if __name__ == '__main__':
    main()
\end{lstlisting}
\end{codefile}

%==============================================================================
\section{Phase 2: LDA模型训练}
%==============================================================================

\subsection{Task 2.1: 构建词典与BOW语料}

\textbf{目标}：为每个时间窗口构建gensim词典，应用词频过滤。

\textbf{输入}：\texttt{data/processed/tokenized/window\_*\_corpus.pkl}

\textbf{输出}：\texttt{output/lda/dictionaries/}
\begin{itemize}
    \item \texttt{window\_*.dict} -- gensim词典
    \item \texttt{window\_*.mm} -- BOW语料（Matrix Market格式）
\end{itemize}

\textbf{状态}：已完成（2026年1月29日）

\textbf{执行结果}：
\begin{itemize}
    \item Phase I (2016-2018): 原始词汇 223,106 → 过滤后 17,137词, 2,131,930 非空文档
    \item Phase II (2021-2022): 原始词汇 379,436 → 过滤后 30,000词, 4,185,913 非空文档
    \item Phase III (2023-2024): 原始词汇 313,353 → 过滤后 27,836词, 2,809,876 非空文档
    \item 过滤参数: min\_df=100, max\_df=0.4, keep\_n=30000
\end{itemize}

\begin{codefile}
\textbf{文件}: \texttt{task2\_modeling/build\_dictionary.py}
\begin{lstlisting}[language=Python]
#!/usr/bin/env python3
# -*- coding: utf-8 -*-
"""
构建gensim词典，应用词频过滤
"""
from gensim import corpora
import pickle
from pathlib import Path

TOKENIZED_DIR = '/Users/yu/code/code2601/TY/Test_LDA/task1_datacleaning/tokenized'
OUTPUT_DIR = '/Users/yu/code/code2601/TY/Test_LDA/task2_modeling/dictionaries'

# 过滤参数（针对百万级数据调优）
MIN_DF = 100      # 最小文档频率（过滤拼写错误/极罕见词）
MAX_DF = 0.4      # 最大文档频率比例（过滤通用词）
KEEP_N = 30000    # 保留词汇上限

def process_window(window_name: str):
    print(f"\n处理 {window_name}...")

    # 加载分词后语料
    corpus_file = Path(TOKENIZED_DIR) / f"{window_name}_corpus.pkl"
    with open(corpus_file, 'rb') as f:
        corpus = pickle.load(f)

    print(f"  文档数: {len(corpus):,}")

    # 构建词典
    dictionary = corpora.Dictionary(corpus)
    print(f"  原始词汇量: {len(dictionary):,}")

    # 过滤极端词频
    dictionary.filter_extremes(
        no_below=MIN_DF,
        no_above=MAX_DF,
        keep_n=KEEP_N
    )
    print(f"  过滤后词汇量: {len(dictionary):,}")

    # 构建BOW语料
    bow_corpus = [dictionary.doc2bow(doc) for doc in corpus]

    # 保存
    Path(OUTPUT_DIR).mkdir(parents=True, exist_ok=True)

    dict_file = Path(OUTPUT_DIR) / f"{window_name}.dict"
    dictionary.save(str(dict_file))

    corpus_file = Path(OUTPUT_DIR) / f"{window_name}.mm"
    corpora.MmCorpus.serialize(str(corpus_file), bow_corpus)

    print(f"  保存: {dict_file}")

    return len(dictionary)

def main():
    windows = ['window_2016_2018', 'window_2021_2022', 'window_2023_2024']

    for window in windows:
        process_window(window)

    print("\n词典构建完成！")

if __name__ == '__main__':
    main()
\end{lstlisting}
\end{codefile}

%------------------------------------------------------------------------------
\subsection{Task 2.2: LDA模型训练}

\begin{note}
训练参数说明：
\begin{itemize}
    \item \texttt{num\_topics=60}: 每个窗口60个主题，平衡粒度与对齐复杂度
    \item \texttt{passes=15}: 迭代次数，确保收敛
    \item \texttt{alpha='auto'}: 自动学习文档-主题分布的稀疏性
    \item \texttt{eta='auto'}: 自动学习主题-词分布的稀疏性
    \item \texttt{chunksize=10000}: 批量大小，适配大数据
\end{itemize}
\end{note}

\begin{codefile}
\textbf{文件}: \texttt{task2\_modeling/train\_lda.py}
\begin{lstlisting}[language=Python]
#!/usr/bin/env python3
# -*- coding: utf-8 -*-
"""
为每个时间窗口训练LDA模型
"""
from gensim import corpora
from gensim.models import LdaMulticore
from gensim.models.coherencemodel import CoherenceModel
import pickle
from pathlib import Path
import logging

logging.basicConfig(level=logging.INFO, format='%(asctime)s - %(message)s')

DICT_DIR = '/Users/yu/code/code2601/TY/Test_LDA/task2_modeling/dictionaries'
TOKENIZED_DIR = '/Users/yu/code/code2601/TY/Test_LDA/task1_datacleaning/tokenized'
OUTPUT_DIR = '/Users/yu/code/code2601/TY/Test_LDA/task2_modeling/models'

# LDA超参数
LDA_CONFIG = {
    'num_topics': 60,
    'passes': 15,
    'chunksize': 10000,
    'alpha': 'auto',
    'eta': 'auto',
    'iterations': 400,
    'random_state': 42,
    'workers': 4,  # 并行数，根据CPU核心调整
}

def train_window(window_name: str):
    print(f"\n{'='*60}")
    print(f"训练 {window_name}")
    print(f"{'='*60}")

    # 加载词典和语料
    dictionary = corpora.Dictionary.load(f"{DICT_DIR}/{window_name}.dict")
    corpus = corpora.MmCorpus(f"{DICT_DIR}/{window_name}.mm")

    print(f"词汇量: {len(dictionary):,}")
    print(f"文档数: {len(corpus):,}")

    # 加载原始分词文本（用于Coherence计算）
    with open(f"{TOKENIZED_DIR}/{window_name}_corpus.pkl", 'rb') as f:
        raw_texts = pickle.load(f)

    # 训练LDA
    print(f"\n开始训练 (num_topics={LDA_CONFIG['num_topics']})...")
    model = LdaMulticore(
        corpus=list(corpus),
        id2word=dictionary,
        **LDA_CONFIG
    )

    # 计算Coherence
    print("\n计算Topic Coherence...")
    coherence_model = CoherenceModel(
        model=model,
        texts=raw_texts[:100000],  # 使用部分样本加速
        dictionary=dictionary,
        coherence='c_v'
    )
    coherence = coherence_model.get_coherence()
    print(f"Coherence (Cv): {coherence:.4f}")

    # 保存模型
    Path(OUTPUT_DIR).mkdir(parents=True, exist_ok=True)
    model_file = f"{OUTPUT_DIR}/{window_name}_lda.model"
    model.save(model_file)
    print(f"模型已保存: {model_file}")

    # 保存Top Words
    topics_file = f"{OUTPUT_DIR}/{window_name}_topics.txt"
    with open(topics_file, 'w', encoding='utf-8') as f:
        f.write(f"Window: {window_name}\n")
        f.write(f"Coherence: {coherence:.4f}\n")
        f.write(f"{'='*60}\n\n")

        for idx in range(model.num_topics):
            top_words = model.show_topic(idx, topn=20)
            words_str = ', '.join([f"{w}({p:.3f})" for w, p in top_words])
            f.write(f"Topic {idx}: {words_str}\n\n")

    print(f"主题词保存: {topics_file}")

    return coherence

def main():
    windows = ['window_2016_2018', 'window_2021_2022', 'window_2023_2024']

    results = {}
    for window in windows:
        coherence = train_window(window)
        results[window] = coherence

    print("\n" + "="*60)
    print("训练结果汇总")
    print("="*60)
    for window, coh in results.items():
        status = "OK" if coh >= 0.40 else "WARN"
        print(f"  {window}: Cv={coh:.4f} [{status}]")

if __name__ == '__main__':
    main()
\end{lstlisting}
\end{codefile}

%==============================================================================
\section{Phase 3: 主题演化对齐}
%==============================================================================

\subsection{Task 3.1: 训练Word2Vec并向量化主题}

\begin{warning}
\textbf{向量空间对齐陷阱}：如果分窗口训练Word2Vec，不同窗口的向量空间坐标轴不同，无法直接计算余弦相似度！

\textbf{解决方案}：使用\textbf{全局Word2Vec}——合并所有窗口语料训练统一模型，确保"Python"在Phase I和Phase III的向量坐标一致。
\end{warning}

\textbf{原理}：将LDA主题的词分布转换为稠密向量，公式：
\[
\vec{V}_{topic} = \sum_{w} P(w|topic) \cdot \vec{V}_{word}
\]

\begin{codefile}
\textbf{文件}: \texttt{task3\_alignment/vectorize\_topics.py}
\begin{lstlisting}[language=Python]
#!/usr/bin/env python3
# -*- coding: utf-8 -*-
"""
主题向量化：LDA主题 -> 稠密向量
"""
import numpy as np
from gensim.models import Word2Vec, LdaMulticore
import pickle
from pathlib import Path

TOKENIZED_DIR = '/Users/yu/code/code2601/TY/Test_LDA/task1_datacleaning/tokenized'
MODEL_DIR = '/Users/yu/code/code2601/TY/Test_LDA/task2_modeling/models'
OUTPUT_DIR = '/Users/yu/code/code2601/TY/Test_LDA/task3_alignment'

# Word2Vec参数
W2V_CONFIG = {
    'vector_size': 200,
    'window': 5,
    'min_count': 50,
    'workers': 4,
    'epochs': 10,
}

def train_word2vec():
    """合并所有语料训练统一的Word2Vec"""
    print("Step 1: 训练Word2Vec...")

    all_corpus = []
    windows = ['window_2016_2018', 'window_2021_2022', 'window_2023_2024']

    for window in windows:
        with open(f"{TOKENIZED_DIR}/{window}_corpus.pkl", 'rb') as f:
            corpus = pickle.load(f)
            all_corpus.extend(corpus)

    print(f"  总文档数: {len(all_corpus):,}")

    model = Word2Vec(sentences=all_corpus, **W2V_CONFIG)

    Path(OUTPUT_DIR).mkdir(parents=True, exist_ok=True)
    model.save(f"{OUTPUT_DIR}/word2vec.model")

    print(f"  词汇量: {len(model.wv):,}")
    print(f"  保存: {OUTPUT_DIR}/word2vec.model")

    return model

def vectorize_topic(lda_model, topic_id: int, w2v_model, topn: int = 50):
    """将单个主题转换为向量"""
    topic_words = lda_model.show_topic(topic_id, topn=topn)

    vector = np.zeros(w2v_model.vector_size)
    total_prob = 0

    for word, prob in topic_words:
        if word in w2v_model.wv:
            vector += prob * w2v_model.wv[word]
            total_prob += prob

    # 归一化
    if total_prob > 0:
        vector = vector / total_prob

    # L2归一化
    norm = np.linalg.norm(vector)
    if norm > 0:
        vector = vector / norm

    return vector

def vectorize_all_windows(w2v_model):
    """向量化所有窗口的主题"""
    print("\nStep 2: 向量化主题...")

    windows = ['window_2016_2018', 'window_2021_2022', 'window_2023_2024']
    topic_vectors = {}

    for window in windows:
        lda_model = LdaMulticore.load(f"{MODEL_DIR}/{window}_lda.model")

        vectors = []
        for topic_id in range(lda_model.num_topics):
            vec = vectorize_topic(lda_model, topic_id, w2v_model)
            vectors.append(vec)

        topic_vectors[window] = np.array(vectors)
        np.save(f"{OUTPUT_DIR}/{window}_topic_vectors.npy", topic_vectors[window])
        print(f"  {window}: shape={topic_vectors[window].shape}")

    # 保存完整字典
    with open(f"{OUTPUT_DIR}/all_topic_vectors.pkl", 'wb') as f:
        pickle.dump(topic_vectors, f)

    return topic_vectors

def main():
    w2v_model = train_word2vec()
    vectorize_all_windows(w2v_model)
    print("\n向量化完成！")

if __name__ == '__main__':
    main()
\end{lstlisting}
\end{codefile}

%------------------------------------------------------------------------------
\subsection{Task 3.2: 混合策略对齐}

\textbf{算法}：
\begin{enumerate}
    \item \textbf{匈牙利算法}：寻找全局最优的1-to-1匹配（Survival）
    \item \textbf{贪婪阈值法}：检测分裂(Split)、融合(Merge)、新生(Birth)、消亡(Death)
\end{enumerate}

\begin{warning}
\textbf{动态阈值}：由于Phase I$\to$II跨度大（隔2年），Phase II$\to$III跨度小（无缝），需使用\textbf{不同阈值}，否则前者会出现大量虚假"消亡"。
\end{warning}

\begin{table}[h]
\centering
\caption{动态对齐阈值设置}
\begin{tabular}{llll}
\toprule
\textbf{转换} & \textbf{Survival阈值} & \textbf{Split/Merge阈值} & \textbf{理由} \\
\midrule
Phase I $\to$ II & 0.55 & 0.45 & 跨度大，允许更大语义漂移 \\
Phase II $\to$ III & 0.70 & 0.60 & 跨度小，要求更严格对应 \\
\bottomrule
\end{tabular}
\end{table}

\begin{codefile}
\textbf{文件}: \texttt{task3\_alignment/align\_topics.py}
\begin{lstlisting}[language=Python]
#!/usr/bin/env python3
# -*- coding: utf-8 -*-
"""
跨时间窗口的主题对齐
算法：匈牙利算法 + 贪婪阈值法
"""
import numpy as np
from scipy.optimize import linear_sum_assignment
from sklearn.metrics.pairwise import cosine_similarity
import pickle
import pandas as pd
from pathlib import Path

VECTOR_DIR = '/Users/yu/code/code2601/TY/Test_LDA/task3_alignment'
OUTPUT_DIR = '/Users/yu/code/code2601/TY/Test_LDA/task3_alignment'

# 动态对齐阈值（根据窗口跨度调整）
THRESHOLDS = {
    ('window_2016_2018', 'window_2021_2022'): {'survival': 0.55, 'split_merge': 0.45},  # 跨度大
    ('window_2021_2022', 'window_2023_2024'): {'survival': 0.70, 'split_merge': 0.60},  # 跨度小
}

def hungarian_alignment(sim_matrix: np.ndarray, threshold: float) -> list:
    """匈牙利算法：1-to-1最优匹配"""
    cost_matrix = 1 - sim_matrix
    row_ind, col_ind = linear_sum_assignment(cost_matrix)

    alignments = []
    for src, tgt in zip(row_ind, col_ind):
        sim = sim_matrix[src, tgt]
        if sim >= threshold:
            alignments.append({
                'source_topic': src,
                'target_topic': tgt,
                'similarity': sim,
                'event': 'survival'
            })

    return alignments

def detect_split_merge(sim_matrix: np.ndarray,
                       survival_alignments: list,
                       threshold: float) -> list:
    """检测分裂/融合/新生/消亡"""
    matched_sources = {a['source_topic'] for a in survival_alignments}
    matched_targets = {a['target_topic'] for a in survival_alignments}

    events = []

    # Split: 已存活的source还有其他高相似target
    for alignment in survival_alignments:
        src = alignment['source_topic']
        primary_tgt = alignment['target_topic']

        for tgt in range(sim_matrix.shape[1]):
            if tgt != primary_tgt and sim_matrix[src, tgt] >= threshold:
                events.append({
                    'source_topic': src,
                    'target_topic': tgt,
                    'similarity': sim_matrix[src, tgt],
                    'event': 'split'
                })

    # Merge: 未匹配的source高度相似于已匹配的target
    for src in range(sim_matrix.shape[0]):
        if src not in matched_sources:
            for tgt in matched_targets:
                if sim_matrix[src, tgt] >= threshold:
                    events.append({
                        'source_topic': src,
                        'target_topic': tgt,
                        'similarity': sim_matrix[src, tgt],
                        'event': 'merge'
                    })

    # Birth: 完全新的target
    for tgt in range(sim_matrix.shape[1]):
        if tgt not in matched_targets:
            max_sim = sim_matrix[:, tgt].max()
            if max_sim < threshold:
                events.append({
                    'source_topic': None,
                    'target_topic': tgt,
                    'similarity': max_sim,
                    'event': 'birth'
                })

    # Death: 完全消失的source
    for src in range(sim_matrix.shape[0]):
        if src not in matched_sources:
            max_sim = sim_matrix[src, :].max()
            if max_sim < threshold:
                events.append({
                    'source_topic': src,
                    'target_topic': None,
                    'similarity': max_sim,
                    'event': 'death'
                })

    return events

def align_windows(vectors_t1, vectors_t2, w1, w2):
    """对齐两个时间窗口（使用动态阈值）"""
    print(f"\n对齐: {w1} -> {w2}")

    # 获取动态阈值
    thresholds = THRESHOLDS.get((w1, w2), {'survival': 0.60, 'split_merge': 0.50})
    survival_th = thresholds['survival']
    split_merge_th = thresholds['split_merge']
    print(f"  阈值: survival={survival_th}, split_merge={split_merge_th}")

    sim_matrix = cosine_similarity(vectors_t1, vectors_t2)

    # Step 1: 匈牙利算法
    survival = hungarian_alignment(sim_matrix, survival_th)

    # Step 2: 分裂/融合检测
    other_events = detect_split_merge(sim_matrix, survival, split_merge_th)

    all_events = survival + other_events

    df = pd.DataFrame(all_events)
    df['source_window'] = w1
    df['target_window'] = w2

    # 统计
    print(f"  Survival: {len([e for e in all_events if e['event']=='survival'])}")
    print(f"  Split: {len([e for e in all_events if e['event']=='split'])}")
    print(f"  Merge: {len([e for e in all_events if e['event']=='merge'])}")
    print(f"  Birth: {len([e for e in all_events if e['event']=='birth'])}")
    print(f"  Death: {len([e for e in all_events if e['event']=='death'])}")

    return df, sim_matrix

def main():
    # 加载向量
    with open(f"{VECTOR_DIR}/all_topic_vectors.pkl", 'rb') as f:
        topic_vectors = pickle.load(f)

    windows = ['window_2016_2018', 'window_2021_2022', 'window_2023_2024']

    all_events = []
    all_matrices = {}

    # 对齐相邻窗口
    for i in range(len(windows) - 1):
        w1, w2 = windows[i], windows[i+1]
        events_df, sim_matrix = align_windows(
            topic_vectors[w1], topic_vectors[w2], w1, w2
        )
        all_events.append(events_df)
        all_matrices[f"{w1}_to_{w2}"] = sim_matrix

    # 保存
    final_df = pd.concat(all_events, ignore_index=True)
    final_df.to_csv(f"{OUTPUT_DIR}/evolution_events.csv", index=False)

    with open(f"{OUTPUT_DIR}/similarity_matrices.pkl", 'wb') as f:
        pickle.dump(all_matrices, f)

    print(f"\n总事件数: {len(final_df)}")
    print(f"保存: {OUTPUT_DIR}/evolution_events.csv")

if __name__ == '__main__':
    main()
\end{lstlisting}
\end{codefile}

%------------------------------------------------------------------------------
\subsection{Task 3.3: 桑基图可视化}

\begin{codefile}
\textbf{文件}: \texttt{task3\_alignment/visualize\_sankey.py}
\begin{lstlisting}[language=Python]
#!/usr/bin/env python3
# -*- coding: utf-8 -*-
"""
绘制技能演化桑基图
"""
from pyecharts import options as opts
from pyecharts.charts import Sankey
import pandas as pd
from gensim.models import LdaMulticore
from pathlib import Path

MODEL_DIR = '/Users/yu/code/code2601/TY/Test_LDA/task2_modeling/models'
ALIGN_DIR = '/Users/yu/code/code2601/TY/Test_LDA/task3_alignment'
OUTPUT_DIR = '/Users/yu/code/code2601/TY/Test_LDA/outputs'

def get_topic_label(lda_model, topic_id: int, topn: int = 3) -> str:
    """获取主题标签（Top 3词）"""
    words = lda_model.show_topic(topic_id, topn=topn)
    return '_'.join([w for w, _ in words])

def main():
    # 加载数据
    events_df = pd.read_csv(f"{ALIGN_DIR}/evolution_events.csv")
    events_df = events_df[events_df['event'].isin(['survival', 'split', 'merge'])]

    # 加载LDA模型
    windows = ['window_2016_2018', 'window_2021_2022', 'window_2023_2024']
    lda_models = {w: LdaMulticore.load(f"{MODEL_DIR}/{w}_lda.model") for w in windows}

    # 构建节点和连接
    nodes = []
    links = []
    node_set = set()

    color_map = {
        'survival': '#5470C6',  # 蓝色
        'split': '#91CC75',     # 绿色
        'merge': '#FAC858',     # 黄色
    }

    for _, row in events_df.iterrows():
        src_w, tgt_w = row['source_window'], row['target_window']

        if pd.notna(row['source_topic']):
            src_label = f"{src_w[:7]}_{get_topic_label(lda_models[src_w], int(row['source_topic']))}"
            if src_label not in node_set:
                nodes.append({'name': src_label})
                node_set.add(src_label)
        else:
            src_label = None

        if pd.notna(row['target_topic']):
            tgt_label = f"{tgt_w[:7]}_{get_topic_label(lda_models[tgt_w], int(row['target_topic']))}"
            if tgt_label not in node_set:
                nodes.append({'name': tgt_label})
                node_set.add(tgt_label)
        else:
            tgt_label = None

        if src_label and tgt_label:
            links.append({
                'source': src_label,
                'target': tgt_label,
                'value': row['similarity'] * 100,
                'lineStyle': {'color': color_map.get(row['event'], '#EE6666')}
            })

    # 生成图表
    Path(OUTPUT_DIR).mkdir(parents=True, exist_ok=True)

    sankey = (
        Sankey(init_opts=opts.InitOpts(width="1800px", height="1000px"))
        .add(
            series_name="",
            nodes=nodes,
            links=links,
            linestyle_opt=opts.LineStyleOpts(opacity=0.5, curve=0.5),
            label_opts=opts.LabelOpts(position="right", font_size=9),
            node_gap=12,
        )
        .set_global_opts(
            title_opts=opts.TitleOpts(title="技能演化桑基图 (2016-2024)"),
            tooltip_opts=opts.TooltipOpts(trigger="item"),
        )
    )

    output_file = f"{OUTPUT_DIR}/skill_evolution_sankey.html"
    sankey.render(output_file)

    print(f"桑基图已生成: {output_file}")
    print(f"  节点数: {len(nodes)}")
    print(f"  连接数: {len(links)}")

if __name__ == '__main__':
    main()
\end{lstlisting}
\end{codefile}

%==============================================================================
\section{Phase 4: 面板数据构建}
%==============================================================================

\subsection{Task 4.1: 统一主题ID映射}

\textbf{目标}：追踪survival链，为跨窗口的"同一技能"赋予统一ID。

\subsection{Task 4.2: 构建面板数据}

\textbf{输出}：\texttt{outputs/panel\_data.csv}

\textbf{字段}：
\begin{itemize}
    \item \texttt{job\_id}: 职位唯一标识
    \item \texttt{window}: 时间窗口
    \item \texttt{unified\_topic\_id}: 统一技能ID
    \item \texttt{topic\_proportion}: 该技能在职位中的比例
    \item \texttt{job\_title}, \texttt{city}, \texttt{salary}, \texttt{education}...
\end{itemize}

%==============================================================================
\section{执行检查清单}
%==============================================================================

\begin{table}[h]
\centering
\caption{执行检查清单}
\begin{tabular}{clll}
\toprule
\textbf{序号} & \textbf{任务} & \textbf{输出文件} & \textbf{状态} \\
\midrule
1.1 & ESCO词典转换（中英双语） & \texttt{data/esco/jieba\_dict/esco\_jieba\_dict.txt} & $\checkmark$ \\
1.2 & 数据按窗口分割 & \texttt{data/processed/windows/window\_*.csv} & $\square$ \\
1.3 & 分词+短语识别 & \texttt{data/processed/tokenized/window\_*\_corpus.pkl} & $\square$ \\
\midrule
2.1 & 构建词典 & \texttt{output/lda/dictionaries/window\_*.dict} & $\square$ \\
2.2 & LDA训练 & \texttt{output/lda/models/window\_*\_lda.model} & $\square$ \\
\midrule
3.1 & Word2Vec+向量化 & \texttt{output/lda/word2vec.model} & $\square$ \\
3.2 & 主题对齐 & \texttt{output/lda/evolution\_events.csv} & $\square$ \\
3.3 & 桑基图 & \texttt{output/lda/skill\_evolution\_sankey.html} & $\square$ \\
\midrule
4.1 & 统一ID映射 & \texttt{output/lda/unified\_topic\_mapping.pkl} & $\square$ \\
4.2 & 面板数据 & \texttt{output/lda/panel\_data.csv} & $\square$ \\
\bottomrule
\end{tabular}
\end{table}

%==============================================================================
\section{综合化程度量化指标}
%==============================================================================

\begin{note}
\textbf{核心假设验证}：如果AI导致岗位技能综合化，Phase III的平均主题熵应显著高于Phase I。
\end{note}

\subsection{指标1：主题熵 (Topic Entropy)}

衡量单个主题内关键词的分布集中度：
\[
H_{topic} = -\sum_{w \in \text{top-}N} P(w|topic) \cdot \log P(w|topic)
\]

\begin{itemize}
    \item 熵\textbf{低} $\rightarrow$ 关键词集中 $\rightarrow$ 专业化主题（如"Java, SQL, 数据库"）
    \item 熵\textbf{高} $\rightarrow$ 关键词分散 $\rightarrow$ 综合化主题（如"Python, 沟通, 项目管理"）
\end{itemize}

\subsection{指标2：跨域共现率 (Cross-Domain Co-occurrence)}

利用ESCO技能分类体系，计算每个主题覆盖的技能大类数量：
\[
\text{CDR}_{topic} = \frac{|\{\text{ESCO大类} : \exists w \in \text{top-}20, w \in \text{该大类}\}|}{|\text{ESCO大类总数}|}
\]

\textbf{示例}：
\begin{itemize}
    \item 2016年"程序员"主题：全是"Java, SQL" $\rightarrow$ CDR = 1/N（仅计算机类）
    \item 2024年"程序员"主题：含"Python, 项目管理, 沟通" $\rightarrow$ CDR = 3/N（计算机+管理+软技能）
\end{itemize}

\subsection{指标3：岗位综合化指数 (Job Hybridization Index)}

基于文档-主题分布$\theta$计算：
\[
\text{JHI}_{doc} = -\sum_{k} \theta_k \cdot \log \theta_k \quad \text{（文档主题熵）}
\]

或使用有效主题数：
\[
\text{ETN}_{doc} = \sum_{k} \mathbf{1}[\theta_k > 0.05]
\]

%==============================================================================
\section{验证标准}
%==============================================================================

\begin{table}[h]
\centering
\caption{质量验证标准}
\begin{tabular}{lll}
\toprule
\textbf{指标} & \textbf{目标值} & \textbf{说明} \\
\midrule
Topic Coherence (Cv) & $\geq 0.40$ & 每个窗口的主题一致性 \\
Survival比例 & 40\%--60\% & 跨窗口存活的主题占比 \\
演化合理性 & 人工审核 & 技术迭代符合行业认知 \\
\bottomrule
\end{tabular}
\end{table}

\begin{table}[h]
\centering
\caption{综合化趋势验证（核心假设）}
\begin{tabular}{lll}
\toprule
\textbf{指标} & \textbf{预期趋势} & \textbf{统计检验} \\
\midrule
平均主题熵 & Phase III $>$ Phase I & Mann-Whitney U \\
平均CDR & Phase III $>$ Phase I & 双样本t检验 \\
平均JHI/ETN & Phase III $>$ Phase I & 配对检验 \\
\bottomrule
\end{tabular}
\end{table}

%==============================================================================
\section{方法论风险点及应对}
%==============================================================================

\subsection{风险点1：主题数量(K)的选择}

\begin{warning}
\textbf{问题}：如果Phase III岗位真的变复杂了，是否需要更多Topic才能覆盖？K在不同窗口不一致会导致熵不可比。
\end{warning}

\textbf{应对策略}：\textbf{固定K作为控制变量}（推荐K=60）。

\textbf{论文表述}：``在控制了主题粒度（K=60）不变的情况下，我们依然观察到每个主题内部的熵显著增加。''——这比动态K更有说服力。

\textbf{补充验证}：可用Coherence Score在每个窗口单独评估K=40/50/60/70/80，确认K=60不是极端选择。

\subsection{风险点2：软技能泛滥的混淆}

\begin{warning}
\textbf{问题}：``综合化''会不会只是因为Phase III出现了更多通用废话（如``沟通能力''、``团队合作''）？这些软技能词增多导致熵增加，而非真正的硬技能跨界。
\end{warning}

\textbf{应对策略}：\textbf{硬/软技能分离分析}

\begin{enumerate}
    \item \textbf{ESCO分类标注}：利用ESCO技能分类体系，将词汇标注为：
    \begin{itemize}
        \item 硬技能（Technical Skills）：编程语言、工具、框架
        \item 软技能（Soft Skills）：沟通、团队、领导力
        \item 领域知识（Domain Knowledge）：行业术语
    \end{itemize}

    \item \textbf{分层熵计算}：
    \[
    H_{hard} = -\sum_{w \in \text{硬技能}} P(w|topic) \log P(w|topic)
    \]
    分别计算硬技能熵和软技能熵。

    \item \textbf{关键检验}：证明\textbf{硬技能熵}在Phase III显著高于Phase I。

    \textbf{论文表述}：``排除软技能后，我们发现硬技能的跨领域共现率（例如：程序员开始要求产品设计能力）在Phase III显著提升。''
\end{enumerate}

\textbf{代码实现要点}：
\begin{itemize}
    \item 在ESCO词典转换时保留技能分类标签
    \item 计算熵时按类别分组
    \item 输出分层结果：总体熵、硬技能熵、软技能熵
\end{itemize}

%==============================================================================
\section{目录结构}
%==============================================================================

\begin{verbatim}
/Users/yu/code/code2601/TY/
+-- config/
|   +-- paths.py                      # 集中路径配置
+-- data/
|   +-- processed/
|   |   +-- cleaned/all_in_one1.csv   # 清洗后数据
|   |   +-- extracted/all_in_one2.csv # 提取后数据
|   |   +-- deduped/all_in_one2_dedup.csv # 去重后数据
|   |   +-- windows/                  # [输出] 分窗口数据
|   |   +-- tokenized/                # [输出] 分词后语料
|   +-- esco/
|       +-- original/                 # ESCO原始数据
|       +-- chinese/                  # 汉化后ESCO
|       +-- jieba_dict/               # jieba词典+元数据
+-- src/
|   +-- cleaning/                     # 数据清洗脚本
|   +-- lda/                          # LDA建模脚本
|   +-- esco/                         # ESCO相关脚本
+-- output/
|   +-- lda/                          # LDA输出
|   +-- esco/                         # ESCO输出
|   +-- reports/                      # 报告
+-- cache/translation/                # 翻译缓存
+-- docs/plan.tex                     # 本文档
+-- scripts/                          # 运行脚本
\end{verbatim}

%==============================================================================
\section{文献参考}
%==============================================================================

\begin{enumerate}
    \item Chun, S. A., et al. (2013). TopicFlow: Visualizing topic alignment of Twitter data over time. \textit{ASONAM}.
    \item El-Assady, M., et al. (2023). TopicRefiner: Coherence-guided steerable LDA. \textit{IEEE TVCG}.
    \item Han, X. (2020). Evolution of research topics in LIS between 1996 and 2019. \textit{Scientometrics}.
    \item Alabdulkareem, A., et al. (2018). Unpacking the polarization of workplace skills. \textit{Science Advances}.
\end{enumerate}

\end{document}
